\section{李群中的形式计算}

令 \(f, g\) 为关于同一组不定元、系数属于 K 的两个形式幂级数，令 \(f_i\)（分别地，\(g_i\)）为 \(f\)（分别地，\(g\)）的 i 次齐次分量。记

\[
f \equiv g \mod \deg p
\]

\[
\text { if } f _ {i} = g _ {i} \text { for } i <   p.
\]

在本节中，G 表示一个 n 维李群芽；设基域 K 的特征为零。我们一次性借助一个坐标图，把 G 中 e 的一个开邻域同 \(K^{n}\) 中 0 的一个开邻域 U 识别，从而 e 识别为 0。对于 U 中的 x, y 以及 \(n \in \mathbf{Z}\)，x.y 表示 x 与 y 的乘积，\(x^{[m]}\) 表示 G 中 x 的 m 次幂（在它们有定义时）。\(x \in U\) 的坐标记为 \(x_{1}, x_{2}, \ldots, x_{n}\)。

